% from ega2-1-fr.tex
\oldpage[II]{105}

  Inversement, si $X$ est projectif sur $Y$, il est quasi-projectif, donc il
  existe une $Y$-immersion $j$ de $X$ dans un $\mathbf{P}_Y^r$
  (\hyperref[II.5.3.3-fr]{5.3.3}) qui est fermée par
  (\hyperref[II.5.4.4-fr]{5.4.4}) et
  (\hyperref[II.5.5.3-fr]{5.5.3}).
  \item[\rm{(iii)}] Le raisonnement de
  (\hyperref[II.5.5.3-fr]{5.5.3}) montre que, pour tout préschéma $Y$, et tout
  entier $r\geq0$, le morphisme structural $\mathbf{P}_Y^r\to Y$ est
  surjectif, car si on pose
  $\mathcal{S}=\mathbf{S}_{\mathcal{O}_Y}(\mathcal{O}_Y^{r+1})$, on a
  évidemment $\mathcal{S}_y=\mathbf{S}_{k(y)}((k(y))^{r+1})$
  (\hyperref[II.1.7.3-fr]{1.7.3}), donc $(\mathcal{S}_n)_y\neq0$ pour tout
  $y\in Y$ et tout $n\geq0$.
  \item[\rm{(iv)}] Il résulte des exemples de Nagata \cite{II-26} qu'il y a
  des morphismes propres qui ne sont pas quasi-projectifs.
\end{enumerate}

\begin{proposition}[5.5.5]
\label{II.5.5.5-fr}
\begin{enumerate}
  \item[\rm{(i)}] Une immersion fermée est un morphisme projectif.
  \item[\rm{(ii)}] Si $f:X\to Y$ et $g:Y\to Z$ sont des morphismes
  projectifs, et si $Z$ est un schéma quasi-compact ou un préschéma dont
  l'espace sous-jacent est noethérien, $g\circ f$ est projectif.
  \item[\rm{(iii)}] Si $f:X\to Y$ est un $S$-morphisme projectif,
  $f_{(S')}:X_{(S')}\to Y_{(S')}$ est projectif pour toute extension
  $S'\to S$ du préschéma de base.
  \item[\rm{(iv)}] Si $f:X\to Y$ et $g:X'\to Y'$ sont des $S$-morphismes
  projectifs, il en est de même de $f\times_Sg$.
  \item[\rm{(v)}] Si $g\circ f$ est un morphisme projectif et si $g$ est
  séparé, $f$ est projectif.
  \item[\rm{(vi)}] Si $f$ est projectif, il en est de même de
  $f_{\mathrm{red}}$.
\end{enumerate}
\end{proposition}

(i) résulte aussitôt de (\hyperref[II.3.1.7-fr]{3.1.7}). Il faut ici
démontrer (iii) et (iv) séparément, à cause de la restriction introduite sur
$Z$ dans (ii) (cf. (\hyperref[I.3.5.1-fr]{I, 3.5.1})). Pour prouver (iii), on
se ramène au cas où $S=Y$ (\hyperref[I.3.3.11-fr]{I, 3.3.11}) et l'assertion
résulte alors aussitôt de (\hyperref[II.5.5.1-fr]{5.5.1, b)}) et de
(\hyperref[II.3.5.3-fr]{3.5.3}). Pour démontrer (iv), on est aussitôt ramené
au cas où $X=\mathbf{P}(\mathcal{E})$, $X'=\mathbf{P}(\mathcal{E}')$,
$\mathcal{E}$ (resp. $\mathcal{E}'$) étant un $\mathcal{O}_Y$-Module (resp.
un $\mathcal{O}_{Y'}$-Module) quasi-cohérent de type fini. Soient $p$, $p'$ les
projections canoniques de $T=Y\times_SY'$ dans $Y$ et $Y'$ respectivement~;
d'après (\hyperref[II.4.1.3.1-fr]{4.1.3.1}), on a
$\mathbf{P}(p^*(\mathcal{E}))=\mathbf{P}(\mathcal{E})\times_YT$,
$\mathbf{P}(p'^*(\mathcal{E}'))=\mathbf{P}(\mathcal{E}')\times_{Y'}T$, d'où
\begin{align*}
\mathbf{P}(p^*(\mathcal{E}))\times_T\mathbf{P}(p'^*(\mathcal{E}'))
&=(\mathbf{P}(\mathcal{E})\times_YT)\times_T
(T\times_{Y'}\mathbf{P}(\mathcal{E}'))\\
&=\mathbf{P}(\mathcal{E})\times_Y
(T\times_{Y'}\mathbf{P}(\mathcal{E}'))
=\mathbf{P}(\mathcal{E})\times_S\mathbf{P}(\mathcal{E}')
\end{align*}
en remplaçant $T$ par $Y\times_SY'$ et utilisant
(\hyperref[I.3.3.9.1-fr]{I, 3.3.9.1}). Or, $p^*(\mathcal{E})$ et
$p'^*(\mathcal{E}')$ sont de type fini sur $T$
(\hyperref[0.5.2.4-fr]{0, 5.2.4}), donc il en est de même de
$p^*(\mathcal{E})\otimes_{\mathcal{O}_T}p'^*(\mathcal{E}')$~; comme
$\mathbf{P}(p^*(\mathcal{E}))\times_T\mathbf{P}(p'^*(\mathcal{E}'))$
s'identifie à un sous-préschéma fermé de
$\mathbf{P}(p^*(\mathcal{E})\otimes_{\mathcal{O}_T}p'^*(\mathcal{E}'))$
(\hyperref[II.4.3.3-fr]{4.3.3}), cela achève de prouver (iv). Pour obtenir (v)
et (vi), on peut appliquer (\hyperref[I.5.5.13-fr]{I, 5.5.13}), car tout
sous-préschéma fermé d'un $Y$-schéma projectif est un $Y$-schéma projectif par
(\hyperref[II.5.5.1-fr]{5.5.1, a)}).

Il reste à prouver (ii)~; en vertu de l'hypothèse sur $Z$, cela résulte de
(\hyperref[II.5.5.3-fr]{5.5.3}), de
(\hyperref[II.5.3.4-fr]{5.3.4, (ii)}) et
(\hyperref[II.5.4.2-fr]{5.4.2, (ii)}).

\begin{proposition}[5.5.6]
\label{II.5.5.6-fr}
Si $X$ et $X'$ sont deux $Y$-schémas projectifs, $X\amalg X'$ est un
$Y$-schéma projectif.
\end{proposition}

C'est une conséquence évidente de (\hyperref[II.5.5.2-fr]{5.5.2}) et
(\hyperref[II.4.3.6-fr]{4.3.6}).

\begin{proposition}[5.5.7]
\label{II.5.5.7-fr}
Soit $X$ un $Y$-schéma projectif, $\mathcal{L}$ un $\mathcal{O}_X$-Module
$Y$-ample~; pour toute section $f$ de $\mathcal{L}$ au-dessus de $X$, $X_f$
est affine sur $Y$.
\end{proposition}

